\section{Characterising the instrument, and searching for what beats the frontier}
\label{sec:instrument}

The first two phases of the v0.7 programme did no policy work at all. One rebuilt the measuring
instrument and established what it can and cannot see; the other tried, in the open, to find
anything that beats the v0.6 \emph{frontier}. Both were run before any candidate architecture
existed, and both returned results that reshaped what the rest of the programme could claim.

\subsection{The frontier, and why the bar is not the deployed policy}
\label{sec:instrument-frontier}

v0.6 shipped a policy with its determinized search switched off, and also \emph{measured} that
search working. The set of configurations available at the end of that cycle is therefore not a
point but a \textbf{frontier}: the deployed blueprint at one end, and progressively more expensive
searching configurations at the other. Three points of it are named throughout this paper --- the
deployed policy \incumbent{}, the cheapest searching configuration on the frontier \fcheap{}, and a
more expensive one \fmid{}.

This distinction is the single most consequential methodological choice in the v0.7 programme.
Beating the deployed policy is the easy claim, and several unfitted one-switch deviations of it
already do. \textbf{The bar was therefore fixed as \fcheap{}}, and \S\ref{sec:protocol-failure}
records that the protocol committed to it in advance.

\subsection{What the instrument can see}
\label{sec:instrument-sees}

The detection-floor apparatus of \S\ref{sec:methodology-floor} was built and calibrated here, and
produced the two findings that section reports: the floor is \vsevenPrimaryFloor{} points, and it
does \emph{not} buy down with evaluation games. Below roughly a point and a half a fitted responder
is empty rather than unresolved, so the binding constraint on every exploitability statement in this
paper is search power.

Two further instrument results matter for how later sections read.

\paragraph{Transcript inversion works, and it is not an attack.} Because the target team is
deterministic and published, an adversary can invert its transcript: for every ask the target made,
ask what deals would have produced that ask, and contract the posterior accordingly. Measured, this
contracts the deal posterior by $\sim$\vsevenBitsPerAsk{} bits per observed ask beyond the
certificates the rules already supply --- the first such measurement in this corpus, and one the
threat model's literature review found no external precedent for. Converting those bits into play is
worth \vsevenInversionUnfitted{}~points unfitted, the strongest single arm anyone built against the
deployed policy.

\textbf{It is not an exploit.} The excess tracks planted-edge size within
$\pm$\vsevenInversionTrack{} on all eight rungs of the calibration ladder, which is the signature of
a general strength gain rather than an exploitation instrument. And no linear handicap adds more
than \vsevenHandicapBits{}~bits per ask, so within this policy family there is \emph{no readability
handicap for a stochastic policy to buy}. A brief that asked for a stochastic policy conditional on
the inverter biting was therefore closed on its antecedent rather than skipped.

\paragraph{A cost figure the corpus had wrong by two orders of magnitude.} The record priced the
searching configuration at \vsevenCorpusCost$\times$ the blueprint. That figure describes the
\emph{unrestricted} search, which is not the point of the frontier anything was measured at. The
endgame-restricted operating point was re-measured at $\sim$\vsevenCostRecorded$\times$.
\S\ref{sec:results-cost} corrects that correction in turn.

\subsection{Open-ended adversary generation}
\label{sec:instrument-adversaries}

The second phase's constraint was that the searches must not share a bias --- \emph{fifteen runs of
one search are one run}. Every exploiter search in the corpus before it had been the same search: a
cross-entropy fit of a linear vector, maximising win rate, from one seed, against one target. The
phase built six adversary classes, three correlation regimes, and a taxonomy of channels, and
measured channel ceilings \emph{before} fitting anything against them.

\paragraph{The headline is negative and it is the reason this paper exists in the shape it does.}
Nothing built exploits the v0.6 frontier beyond the detection floor. The best measured exploitability
of the deployed policy remains an in-class figure of \vsevenBestExploit{}
$[\vsevenBestExploitLo, \vsevenBestExploitHi]$ over \vsevenBestExploitN{} games --- \emph{below} the
class floor of \vsevenPrimaryFloor{}. Seven independently fitted searches confirmed it; none cleared
its class floor.

Every arm that beat the deployed policy by more than the floor turned out, on measurement, to be a
\textbf{better policy rather than an exploit}: the target's own test-time search at
\vsevenSearchAsExploit{}, and a single out-of-class coordinate at \vsevenOneSwitchArm{} whose gain
survives against opponents sharing none of the target's machinery. The second of those is worth
dwelling on, because it is where this cycle's entire strength gain comes from: it is a
\emph{hand-set} coordinate found by an unfitted sweep, with no fit at all, and against the more
expensive frontier point its edge is \vsevenOneSwitchArmVsFmid{} --- negative. It is a strength gain
with a worst case, not an exploit.

\paragraph{A defect found in the target rather than an attack on it.} The declaration-urgency channel
the phase expected to headline has a measured ceiling of \vsevenUrgencyCeiling{}
$[\vsevenUrgencyCeilingLo, \vsevenUrgencyCeilingHi]$ --- dead as an attack. But the \emph{same}
channel is worth \vsevenUrgencyWorthLo{} to \vsevenUrgencyWorthHi{} points to a configuration that
simply switches it off. The adversary owns the smaller half of a two-point defect in the target.

\paragraph{A cliff nobody can reach, and a termination duty inherited by anyone who removes it.}
Behind the incumbent's event counter sits a \vsevenCliff{}-point cliff: past a threshold its
declaration rule drops its team-ownership floor and cashes half-suits at close to a coin flip. No
adversary reaches it --- the longest game any of them produced against an unmodified incumbent is
\vsevenCliffIncumbentTail{} events. But that escalation is also the policy's termination guarantee,
and a configuration that switches it off has a self-play tail of \vsevenCliffTail{} events. Removing
a defect here creates an obligation, and \S\ref{sec:configuration-keys} is how it was discharged.

\subsection{What these two phases established before any candidate was built}
\label{sec:instrument-established}

\begin{enumerate}
\item The v0.7 case rests on \textbf{beating the frontier, not on closing exploits}. There are no
      exploits above the floor to close.
\item The bar is \fcheap{}, fixed in advance.
\item The only thing found that beats the frontier is a configuration built from the frontier's own
      search plus one hand-set coordinate plus a switch --- a \emph{configuration change}, not an
      architecture. This is the object called \composite{} throughout this paper, and
      \S\ref{sec:against-composite} is the section in which it becomes decisive.
\item The instrument's floor is \vsevenPrimaryFloor{} points and will not improve by spending games.
\item The evaluation material is sealed, and the adversary population is split in half by a rule
      fixed in advance, with half of it encrypted (\S\ref{sec:methodology-seal}).
\end{enumerate}
